\clearpage
\phantomsection

\addcontentsline{toc}{chapter}{Tóm tắt}
\chapter*{\fontsize{13}{13}\selectfont{Tóm tắt}}

Trí tuệ nhân tạo (AI) đang trở thành xu hướng phát triển mạnh mẽ, trong đó việc ứng dụng Học tăng cường sâu (Deep Reinforcement Learning) để xây dựng các tác nhân thông minh đóng vai trò quan trọng trong các bài toán ra quyết định tuần tự. So với các phương pháp lập trình kịch bản truyền thống, tác nhân học tăng cường có khả năng tự cải thiện thông qua tương tác với môi trường và thích nghi tốt hơn trước các tình huống không được lập trình sẵn. Khóa luận này nghiên cứu bài toán huấn luyện CombatAgent trong một game hành động Roguelike góc nhìn từ trên xuống được xây dựng bằng Unity, TopDown Engine và Unity ML-Agents. Tác nhân sử dụng quan sát vector 207 chiều, không gian hành động rời rạc nhiều nhánh $(3,3,3,9)$ và được huấn luyện bằng PPO.

Điểm chính của khóa luận là kết hợp PCGNN với curriculum learning. PCGNN được dùng để sinh bản đồ offline, sau đó adapter chuyển kết quả sinh sang định dạng văn bản của game, kiểm tra khả năng đi tới bằng BFS và đánh giá độ khó bằng các chỉ số như tỷ lệ tường, độ dài đường đi ngắn nhất, tỷ lệ ngõ cụt và A* difficulty. Các bản đồ hợp lệ được chia thành Easy, Medium và Hard để đưa vào quá trình học theo chương trình trong Unity. Kết quả thực nghiệm trên các run ablation cho thấy curriculum learning là yếu tố tạo cải thiện lớn nhất: PPO baseline đạt reward tail-50 là 2.073, trong khi run curriculum đạt Hard mean 3.843, tương đương mức cải thiện khoảng 85\%. Khi kết hợp ICM trên nền curriculum, reward Hard tăng lên 4.369, cho thấy phần thưởng nội tại có thể hỗ trợ thêm khi tác nhân đã được scaffold bằng các lesson dễ hơn. Ngoài ra, khóa luận ghi nhận hiện tượng entropy plateau quanh 4.40 nat, gợi ý rằng thiết kế action space và cấu trúc policy là hướng cần nghiên cứu tiếp.

\vspace{0.5cm}
\noindent{\textbf{Từ khóa:}} \textit{Học tăng cường sâu, Game Roguelike, Unity ML-Agents, PPO, PCGNN, Procedural Content Generation, Curriculum Learning}
